\documentclass[a4paper,12pt]{article}
\usepackage{amsmath, amsthm, amssymb, geometry}
\usepackage{xcolor}
\usepackage{float}
\usepackage[pdfinfo=on]{xepersian}
\settextfont{XB Niloofar}
\setlatintextfont{Times New Roman}

\geometry{margin=2.5cm}
\linespread{1.2}

\newtheoremstyle{exercise}{3pt}{3pt}{\normalfont}{}{\bfseries}{\\}{0pt}{}
\theoremstyle{exercise}
\newtheorem{exercise}{تمرین}
\newtheorem{solution}{پاسخ}

% غیرفعال کردن شماره‌گذاری
\renewcommand{\theexercise}{}
\renewcommand{\thesolution}{}

\title{%
	\Huge \textbf{تمرین} \\[0.3cm]
	\Large مبانی ریاضیات - منطق گزاره‌ها
}
\author{علی عاشوری}
\date{\today}

\begin{document}
	
	\maketitle
	
	\section*{وضعیت}
	\begin{itemize}
		\item \textbf{درس:} مبانی ریاضیات
		\item \textbf{فصل:} منطق گزاره‌ها
		\item \textbf{شماره‌ی تمرین:} ۱
		\item \textbf{تاریخ حل:} \lr{July 8, 2026}
		\item \textbf{وضعیت:} \textcolor{blue}{حل شده}
		\item \textbf{منبع:} جزوه‌ی مبانی ریاضیات، دکتر محسن خانی، تمرین ۱.۲.
	\end{itemize}
	\rule{\textwidth}{0.4pt}
	
	\begin{exercise}
		جدول ارزش هر یک از گزاره‌های زیر را رسم کنید.
		\begin{enumerate}
			\item $((p \rightarrow q) \rightarrow (\neg(q \rightarrow p)))$
			\item $((p \rightarrow q) \rightarrow q)$
			\item $((p \rightarrow q) \wedge (p \rightarrow (\neg q)))$
			\item $(\neg (p \rightarrow q))$
			\item $(\bot \rightarrow p)$
			\item $(p \rightarrow \bot)$
		\end{enumerate}
	\end{exercise}
	
	\begin{solution}
			۱.
		\begin{latin}
			\begin{table}[H]
				\centering
				\begin{tabular}{c|c|c|c|c|c}
					$p$ & $q$ & $(p \rightarrow q)$ & $(q \rightarrow p)$ & $(\neg(q \rightarrow p))$ & $((p \rightarrow q) \rightarrow (\neg(q \rightarrow p)))$ \\
					\hline
					T & T & T & T & F & F \\
					T & F & F & T & F & T \\
					F & T & T & F & T & T \\
					F & F & T & T & F & F \\
				\end{tabular}
			\end{table}
		\end{latin}
		
		\noindent
		۲.
		\begin{latin}
			\begin{table}[H]
				\centering
				\begin{tabular}{c|c|c|c}
					$p$ & $q$ & $(p \rightarrow q)$ & $((p \rightarrow q) \rightarrow q)$ \\
					\hline
					T & T & T & T \\
					T & F & F & T \\
					F & T & T & T \\
					F & F & T & F \\
				\end{tabular}
			\end{table}
		\end{latin}
		
		\noindent
		۳.
		\begin{latin}
			\begin{table}[H]
				\centering
				\begin{tabular}{c|c|c|c|c|c}
					$p$ & $q$ & $(\neg q)$ & $(p \rightarrow q)$ & $(p \rightarrow (\neg q))$ & $((p \rightarrow q) \wedge (p \rightarrow (\neg q)))$ \\
					\hline
					T & T & F & T & F & F \\
					T & F & T & F & T & F \\
					F & T & F & T & T & T \\
					F & F & T & T & T & T \\
				\end{tabular}
			\end{table}
		\end{latin}
		
		\noindent
		۴.
		\begin{latin}
			\begin{table}[H]
				\centering
				\begin{tabular}{c|c|c|c}
					$p$ & $q$ & $(p \rightarrow q)$ & $(\neg (p \rightarrow q))$ \\
					\hline
					T & T & T & F \\
					T & F & F & T \\
					F & T & T & F \\
					F & F & T & F \\
				\end{tabular}
			\end{table}
		\end{latin}
		
		\noindent
		۵.
		\begin{latin}
			\begin{table}[H]
				\centering
				\begin{tabular}{c|c|c}
					$\bot$ & $p$ & $(\bot \rightarrow p)$ \\
					\hline
					F & T & T \\
					F & F & T
				\end{tabular}
			\end{table}
		\end{latin}
		
		\noindent
		۶.
		\begin{latin}
			\begin{table}[H]
				\centering
				\begin{tabular}{c|c|c}
					$p$ & $\bot$ & $(p \rightarrow \bot)$ \\
					\hline
					T & F & F \\
					F & F & T
				\end{tabular}
			\end{table}
		\end{latin}
	\end{solution}
	
	\section*{یادداشت‌ها}
	\begin{itemize}
		\item گزاره‌ی ۱ یعنی $((p \rightarrow q) \rightarrow (\neg(q \rightarrow p)))$ فقط در صورتی درست است که فقط یکی از گزاره‌های p یا q درست باشد. به عبارت دیگر، این گزاره معادل یای مانع جمع است و با گزاره‌ی زیر معادل است:
		$$ (p \vee q) \wedge (\neg (p \wedge q)) $$
	\end{itemize}
	
\end{document}
